%% notes-tex/030-scaling-study/sections/4-support.tex

\section{What each fit can support\secstatus{todo}}
\label{sec:support}

Three defects would have stopped every figure, and each was reproduced before it
was fixed. The submit script passed \texttt{transforms.fit\_stats\_strict=false}
without the \texttt{+} Hydra requires for a key no config defines; composing the
exact launch overrides locally reproduced the refusal, and the first calibration
job (22812512) printed \texttt{Could not override 'transforms.fit\_stats\_strict'}
after thirty minutes of imports from Lustre. The config now defines the key, and
since configs and code are read from the Delta checkout when a job starts, a
pull repaired every pending job in place. The compute callback logged its keys
without an epoch field, and a W\&B history query for \texttt{epoch} together with
\texttt{compute/records\_seen} returned zero rows on the one instrumented run
(53qqijnu) while the key alone returned 28; the callback now logs the epoch on
every row and the readout falls back to the epochs-completed counter for older
runs. The trainer logged a per-token Pearson but no per-token MSE, so the
per-screen keys the readout fits did not exist, and the first missing key
crashed the reads on \texttt{float(None)}; the per-token collection now carries
MSE in label units beside Pearson. A fourth, the unreachable matched budget, is
above. None of these touched the model or the training.

With those repaired, the three fits divide by how much the queued runs support.

\notesec{The paired gap: supported at 1/64 and 1/16}
Both builds have three seeds at the two smallest subsets, both converge there,
and the reads are at the same $B$. The gap is interpretable at those two points,
with a seed spread on each side. At 1/4 it is one seed each. At the full pool it
is the difference of two unconverged ten-epoch runs, a gap in learning speed at
ten epochs rather than in attainable loss, and it is drawn hollow. Two confounds
belong in its caption: the 025 normalizer refits on each subset while 030 pins
the full-pool constants, and the 030 arm differs from 025 by more than the
per-entry rows, since it also carries the dataset token and the Kuzmin
query-strain double-fitness entries the 025 build never had. The gap is
``per-entry rows plus token plus the added measurement types''; separating the
token is one more 180x8 column.

\notesec{$L(D)$ at 180x8: supported on the matched read, in a narrower claim}
Four points at 180x8, with three seeds at the two smallest. On the selected read
the small subsets are converged and the full pool is not, so the curve flattens
at its top and $\beta$ is biased low with the slack absorbed into $E$; the
minimum over 240 noisy epochs against a minimum over ten biases the same way. On
the matched read every point is a loss at $B = 5.12$M records seen, and the fit is
the claim it can be: at a fixed budget, what more distinct data buys. The deeper
limit is identifiability. The pooled replicate floor works out near 1.7e-3 to
1.9e-3 in label units; the best pooled validation MSE this arm has reached is
3.3e-3 (53qqijnu, epoch 40), and predicting the per-screen mean scores about
3.9e-3. Loss minus floor therefore spans less than a factor of 2.5 across the
whole curve, and the floored form has one degree of freedom on four points. The
power form's held-out test is the only one that means anything on this axis, and
the fitted $E$ is a bound to report with its interval rather than a value.

\notesec{$L(N)$ at the full pool: a fixed-epoch slice}
Five sizes, every one at ten epochs, because the budget table is keyed by subset.
If the GPU is the bottleneck, the small sizes finish ten epochs in a fraction of
48 h and could have run to convergence; if the loader is, as it was on IGB's
GPFS, every size costs 4.3 h an epoch and ten epochs is what any of them gets.
The two calibration jobs tell which. Either way the slice measures loss after
10.5M records seen as a function of $N$, which conflates how fast a size learns
per epoch with how much it can learn, and the learning rate is checked at the
smallest size only while every size runs at a batch four times smaller than the
one the arm was tuned at. The exponent from this slice is descriptive until the
small sizes have been trained to convergence and the learning rate has been
checked at the top of the ladder.

\notesec{The compute axis: not in this grid}
IsoFLOP curves need several budgets with several sizes at each, with records
seen set by the budget and the per-size cost per record from the calibration
jobs. The grid logs everything such a sweep would read, and the fixed-epoch
column is one budget line of it, but a frontier needs three.
